% !TeX root = ../教材题库_全部条目_转写工作稿.tex
% Source: XJ-B1-C06-EX-P236-Y0767
% 第6章 / 6.4 用样本估计总体 / 教材 p.236 / PDF 第243页
% OCR未校对：公式与文字均可能有误；未完成数学TeX转写；禁止直接发给学生
% 原题及相邻原解截图（核对用）：../题目原图/XJ-B1-C06-EX-P236-Y0767-01.webp
% 本文件只在转写完成后作为题干使用；OCR 条目尚不是可靠题干。
\begin{bankitem}{6.4 例题 例题1}{第6章 / 6.4 用样本估计总体 / 教材 p.236 / PDF 第243页}{XJ-B1-C06-EX-P236-Y0767}{OCR未校对：公式与文字均可能有误；未完成数学TeX转写；禁止直接发给学生}
\begin{ocrdraft}
[教材p.236 / PDF第243页]
\par
0〔〔]一种产品需要人工组装，现有A，B两种可供选择的组装方法．为检验
\par
哪种方法生产效率更高，现随机抽取29名工人并随机分成两组：第一组14人，采
\par
用方法A组装；第二组15人，采用方法B组装．让两组工人在相同的时间内组装
\par
产品，得到组装好的产品数量（单位：个）如下表所示：
\par
方法126129129130131127129127128128127128128125
\par
方法B129125126126119126128127126127127126126127125
\par
哪种组装方法的效率更高？
\par
分析平均数刻画了一组数据的平均水平．我们要比较组装方法在相同时间内
\par
的效率时，可以分别计算用不同组装方法得到的产品数量的平均数，再通过平均数
\par
来进行比较．
\par
解设两组工人采用方法，A，B组装的平均产量分别为云，到，
\par
126十129十129十，．十125
\par
一128（个），
\par
14
\par
129十125十126十“，十125
\par
一126（个）．
\par
由于在相同时间内，方法，A的平均产量高于方法B的平均产量，因而可以认为
\par
方法A的效率更高\ensuremath{\cdot}
\end{ocrdraft}
\end{bankitem}
